\documentclass[10pt,a4paper]{article}
\usepackage[margin=16mm]{geometry}
\usepackage[T1]{fontenc}
\usepackage{lmodern}
\usepackage[ngerman]{babel}
\usepackage{booktabs,tabularx,array}
\usepackage{xcolor}
\usepackage{hyperref}
\setlength{\parindent}{0pt}
\setlength{\parskip}{4pt}
\renewcommand{\arraystretch}{1.15}

\begin{document}
{\Large\textbf{Bump'em: Anforderungen an die Stromversorgung}}\\[2pt]
IBMS-Labor, Hochschule Offenburg \quad$\cdot$\quad für die Elektrofachkraft \quad$\cdot$\quad 21.~Oktober 2026 \quad$\cdot$\quad Lorenzo Lucchini

\section*{1. Das System}
Vier baugleiche Türme (A, B, C, D) stehen um das GRAIL-Laufband. Jeder Turm zieht mit einem Elektromotor an einem Seil, das an
einer Person auf dem Laufband befestigt ist. Netzseitige Verbraucher sind nur die unten aufgeführten Netzteile; alle anderen
Komponenten werden über deren Gleichspannungsausgänge versorgt (48\,V Motortreiber, 24\,V Kraftmessverstärker). Ein Steuer-PC
(Laptop mit eigenem Netzteil) ist per USB mit den Türmen verbunden.

\section*{2. Netzseitige Verbraucher je Turm (Herstellerangaben)}
\begin{tabularx}{\linewidth}{>{\raggedright\arraybackslash}p{4.2cm}>{\raggedright\arraybackslash}X>{\raggedright\arraybackslash}X}
\toprule
 & \textbf{Mean Well RSP-2000-48} (48-V-Motornetzteil) & \textbf{Mean Well RS-15-24} (24-V-Netzteil für den Messverstärker) \\
\midrule
Ausgangsnennwerte & 48\,V, 42\,A, 2016\,W & 24\,V, 0,625\,A, 15\,W \\
Eingangsspannung / Frequenz & 90--264\,VAC (volle Leistung ab 200\,VAC), 47--63\,Hz & 85--264\,VAC, 47--63\,Hz \\
Eingangsstrom laut Typenschild & \textbf{12,3\,A} bei 200--240\,VAC & 0,35\,A bei 100--240\,VAC \\
Eingangsstrom typ.\ bei Volllast & 10\,A bei 230\,VAC & 0,25\,A bei 230\,VAC \\
Einschaltstrom (Kaltstart, typ.) & \textbf{50\,A} & \textbf{65\,A} bei 230\,VAC \\
Ableitstrom & $<$ 2\,mA bei 240\,VAC & $<$ 2\,mA bei 240\,VAC \\
Leistungsfaktor / Wirkungsgrad & 0,97 bei 230\,VAC Volllast / 92\,\% & -- / 82\,\% \\
Netzanschluss & Schraubklemmen L, N, PE (Schutzklasse~I) & Schraubklemmen L, N, PE; \textbf{offen, nicht abgedeckt} \\
Sonstiges & interne Lüfter; Überlastschutz 105--125\,\% & Leerlaufleistung $<$ 0,5\,W \\
\bottomrule
\end{tabularx}

\smallskip
\textbf{Je Turm, laut Typenschild:} 12,3 + 0,35 $\approx$ \textbf{12,7\,A} bei 230\,V. \quad
\textbf{Alle vier Türme, laut Typenschild:} $\approx$ \textbf{51\,A}.\\
Quellen: Mean-Well-Datenblätter RSP-2000 und RS-15 sowie die Typenschilder der Geräte.

\section*{3. Tatsächliche Leistungsaufnahme dieses Systems}
Das Typenschild gibt an, was die Netzteile \emph{liefern könnten}. In diesem System ist der Motorstrom durch die Konfiguration
der Motortreiber auf 9,9\,A je Motor begrenzt; das begrenzt die Leistungsaufnahme jedes Turms:

\begin{tabularx}{\linewidth}{>{\raggedright\arraybackslash}X r r}
\toprule
Betriebszustand & je Turm & 4 Türme gleichzeitig \\
\midrule
Eingeschaltet, Leerlauf (kein Zug; von Mean Well nicht angegeben) & gering & gering \\
Seil leicht gespannt halten (Normalbetrieb) & $<$ 50\,W & $<$ 200\,W \\
\textbf{Maximaler Zug, aktuelle Einstellung} (Motor 9,9\,A bei 48\,V) & $\le$ 520\,W ($\approx$ 2,3\,A) & $\le$ \textbf{2,1\,kW ($\approx$ 9,3\,A)} \\
Mögliche spätere Einstellung (Software-Kraftgrenze 200\,N, $\approx$ 16,8\,A Motorstrom) & $\le$ 870\,W ($\approx$ 3,8\,A) & $\le$ 3,5\,kW ($\approx$ 15\,A) \\
\bottomrule
\end{tabularx}

\smallskip
Maximale Züge sind kurz: Ein Zug dauert höchstens 0,6\,s, typischerweise einige pro Minute. Die Bremsenergie des Motors wird
auf der 48-V-Seite in einem Widerstand (Bremschopper) in Wärme umgesetzt; es wird nichts ins Netz zurückgespeist. Die Werte
sind unsere Abschätzung aus den Antriebseinstellungen, keine Messung; bei Bedarf messen wir mit einem
Zwischenstecker-Leistungsmessgerät nach.

\smallskip
\small Anordnung der Türme: A links, B vorne, C rechts, D hinten am Laufband. Datenblätter und die vollständige technische
Dokumentation stellen wir gerne zur Verfügung.
\end{document}
